\chapter{Introducción}

La creciente preocupación por el cambio climático y la necesidad de reducir las emisiones de gases de efecto invernadero han impulsado una transición global hacia sistemas energéticos basados en fuentes renovables. En este contexto, tecnologías como la energía solar fotovoltaica y la energía eólica han adquirido un rol protagónico debido a su disponibilidad, sostenibilidad y rápida expansión a nivel mundial.

No obstante, la naturaleza intermitente y variable de estos recursos plantea desafíos significativos en términos de estabilidad, confiabilidad y gestión de la energía. Para abordar estas limitaciones han surgido los sistemas híbridos de energías renovables, que integran múltiples fuentes de generación junto con sistemas de almacenamiento. Entre las alternativas disponibles, el hidrógeno verde ha ganado especial interés como vector energético, debido a su capacidad de almacenamiento estacional, su alta densidad energética másica y su versatilidad de uso, que permite descarbonizar sectores en los que la electrificación directa resulta inviable \cite{Irena2020GreenHydrogen}.

Los sistemas híbridos que combinan generación eólica, solar, almacenamiento electroquímico y tecnologías de hidrógeno permiten mejorar la flexibilidad operativa y el aprovechamiento del recurso renovable \cite{Wang2021HybridSystem,Naderipour2021HybridClean}. Sin embargo, su dimensionamiento óptimo constituye un problema complejo, debido a la interacción entre múltiples componentes, la variabilidad temporal de los recursos y la presencia de restricciones técnicas, económicas y regulatorias acopladas entre sí.

En este contexto, las metaheurísticas constituyen la familia de métodos habitualmente empleada para abordar problemas no lineales y no convexos como el dimensionamiento de sistemas HRES-H$_2$ \cite{Khan2022Review}. Sin embargo, una comparación limitada a valores medios o mejores ejecuciones no permite separar desempeño sistemático de variación estocástica, y la ausencia de un óptimo conocido impide saber cuán lejos quedan las soluciones halladas. Esta tesis incorpora una línea base aleatoria, múltiples ejecuciones independientes por algoritmo y el óptimo exacto de cada instancia, obtenido por enumeración completa del dominio discretizado, para delimitar lo que los datos permiten afirmar.

Esta carencia adquiere relevancia particular a la luz del teorema de la no existencia de almuerzo gratis \cite{Wolpert1997NFL}, según el cual ninguna metaheurística es universalmente superior. Si la superioridad de un algoritmo sólo puede establecerse empíricamente y para una clase acotada de problemas, entonces el rigor del protocolo experimental no es un detalle accesorio del trabajo: es la condición que determina si sus conclusiones son válidas. La misma lógica se aplica a la hibridación de metaheurísticas \cite{Talbi2002Taxonomy}: combinar dos métodos no garantiza mejorar a ninguno, y sólo una comparación del híbrido frente al método del que parte, bajo el mismo presupuesto, permite atribuir un efecto al componente añadido.

\section{Contribuciones del trabajo}

Esta tesis aborda el dimensionamiento de un sistema híbrido eólico-fotovoltaico con electrólisis y almacenamiento electroquímico, y realiza cinco contribuciones:

\begin{enumerate}
    \item \textbf{Una reproducción parcial y trazable} del modelo WPEB de Li et al. \cite{Li2024WPEB}, con errores absolutos y relativos, comparación de los horizontes de 8.760 y 8.784 horas y declaración explícita de la incertidumbre introducida por las coordenadas y detalles de despacho no publicados.

    \item \textbf{Una formulación mixta continua--discreta del problema y cuatro instancias} en las que la capacidad eólica permanece continua y el electrolizador, la potencia de batería y su duración pertenecen a catálogos discretos. Las instancias varían el límite de ventas a red, la finura de los catálogos y la escala de la planta, y cada una se resuelve por fuerza bruta sobre el dominio discretizado (\BruteForceDesignsTotal{} diseños anuales en total), de modo que el óptimo exacto de cada instancia es el referente de todas las brechas.

    \item \textbf{Tres metaheurísticas híbridas} construidas sobre una infraestructura común de hibridación: un algoritmo genético memético (GA+LS), un enjambre de partículas con operadores genéticos (PSO+GA) y un optimizador de hienas con búsqueda local (SHO+LS). Cada híbrido conserva intacto el método del que parte y respeta exactamente el mismo presupuesto de solicitudes, lo que hace atribuible al componente añadido cualquier diferencia observada.

    \item \textbf{Un banco de pruebas comparativo} entre los \NAlgorithms{} algoritmos sobre las \NInstances{} instancias, con \RequestsPerRun{} solicitudes y \Replicates{} ejecuciones independientes por método e instancia (\TotalDesignRequests{} diseños anuales evaluados). El protocolo compara mediana, rango intercuartílico, brecha y tasa de éxito respecto del óptimo exacto, y una familia planificada de contrastes: cada metaheurística frente a la línea base y cada híbrido frente a su padre.

    \item \textbf{Un análisis de sensibilidad} temporal y de modelado que evita presentar un ordenamiento observado como propiedad universal de los algoritmos.
\end{enumerate}

\section{Estructura del documento}

El Capítulo 2 formula el problema de dimensionamiento, define sus cuatro instancias y caracteriza su dificultad computacional. El Capítulo 3 enuncia el objetivo general y los objetivos específicos. El Capítulo 4 desarrolla el marco teórico en tres bloques: los fundamentos del hidrógeno verde y de los sistemas HRES-H$_2$, la formulación matemática completa del modelo WPEB, y los fundamentos de la optimización metaheurística, incluida la taxonomía de los híbridos, junto con las brechas de investigación identificadas.

El Capítulo 5 describe la metodología, el entorno computacional, las fuentes de datos, las decisiones de modelado y la arquitectura de software que hace comparables a los nueve optimizadores. El Capítulo 6 define el diseño experimental: las instancias, el óptimo exacto por fuerza bruta, la configuración de los algoritmos clásicos e híbridos y los protocolos experimental y estadístico. El Capítulo 7 presenta y discute los resultados por instancia y de forma transversal. El Capítulo 8 recoge las conclusiones, el aporte del trabajo a la Ingeniería Informática, sus limitaciones y las líneas de trabajo futuro.
